\documentclass[main.tex]{subfiles}
\begin{document}
% ============================================================
%  Appendix: spectral arguments for mesh independence
% ============================================================
\section{Appendix: Spectral Arguments}

\begin{frame}[label=spectral]{Discrete Hessians: A Uniform Spectral Interval}
  \begin{hlbox}
    \textbf{Stability of Galerkin.} Let $V_h \subset \Hzero$ and $u_h = S_h z \in V_h$: \; $a(u_h, v_h) = (z, v_h)_{L^2}$ for all $v_h \in V_h$. Test with $v_h = u_h$:
    \[
      \norm{\nabla u_h}_{L^2}^2 = (z, u_h)_{L^2} \le \norm{z}_{L^2}\norm{u_h}_{L^2} \le c_F\norm{z}_{L^2}\norm{\nabla u_h}_{L^2}
    \]
    so $\norm{\nabla u_h}_{L^2} \le c_F\norm{z}_{L^2}$ and $\norm{u_h}_{L^2} \le c_F\norm{\nabla u_h}_{L^2} \le c_F^2\,\norm{z}_{L^2}$.
  \end{hlbox}

  \visible<2->{
  \begin{hlbox}
    Conforming Galerkin \emph{inherits} coercivity: $\norm{CS_h}_{\mathcal{L}(L^2)} \le c_F^2$, independent of $h$.
    So $H_h := \Lambda^{-1}(CS_h)^*\Lambda(CS_h) + \alpha I$ has $\sigma(H_h) \subset [\alpha,\, \alpha + c_F^4]$ for every mesh (for P1 controls, its matrix is $M^{-1}S_h^\top M S_h + \alpha I$).
  \end{hlbox}
  }

  \visible<3->{
  \begin{qbox}
    Hence $q_h \le \max\{|1 - s\alpha|, |1 - s(\alpha + c_F^4)|\}$ for \emph{all} meshes.
  \end{qbox}
  }
  \vfill\hfill\hyperlink{disch<3>}{\beamerreturnbutton{Back to Section 6}}
\end{frame}

\begin{frame}[label=euclidapp]{Euclidean Hessian: The Derivation}
  \begin{hlbox}
    \textbf{Euclidean Hessian.} $H_E = S_h^\top M S_h + \alpha M = M H_h$, where $H_h$ represents the Hessian in the $M$-inner product. By the Rayleigh quotient bound for $H_h$ ($|\cdot|$ the Euclidean norm),
    \[
      x^\top H_E x = (x, H_h x)_M \in [\alpha,\, \alpha + c_F^4]\cdot x^\top M x \subset \big[\alpha\lambda_{\min}(M),\, (\alpha + c_F^4)\lambda_{\max}(M)\big]\,|x|^2 .
    \]
  \end{hlbox}

  \begin{hlbox}
    Hence $\mathrm{cond}(H_E) \le (1 + c_F^4/\alpha)\,\mathrm{cond}(M)$. On quasi-uniform meshes $\mathrm{cond}(M)$ is bounded (only the $h^d$ scaling hurts); on graded meshes $\mathrm{cond}(M) \sim (h_{\max}/h_{\min})^d$. The factor is attained asymptotically: for oscillatory $x$, $\norm{CS_h x}_{L^2}^2 \ll \alpha\, x^\top M x$, so $H_E \approx \alpha M$ there.
  \end{hlbox}
  \vfill\hfill\hyperlink{disch<3>}{\beamerreturnbutton{Back to Section 6}}
\end{frame}

\end{document}
